\documentclass[10pt]{article}
\usepackage[T1]{fontenc}
\usepackage[utf8]{inputenc}
\usepackage[french]{babel}
\usepackage{amsmath,amssymb,mathrsfs,mathpazo}
\usepackage{geometry}
\geometry{a4paper,margin=2cm}
\usepackage{graphicx,array,tabularx,tabularray,colortbl,multirow,multicol,booktabs,enumitem,pifont,ulem,fancybox,diagbox,cancel}
\usepackage{pstricks,pstricks-add,pst-plot,pst-tree,pst-node,pst-text,pst-func,pst-math,pst-bspline,pst-3dplot,pst-eucl,pst-solides3d}
\let\pscalculate\undefined
\usepackage{pst-bezier}
\usepackage{tikz,pgfplots,tkz-tab}
\usetikzlibrary{arrows,arrows.meta,patterns,calc,positioning,trees,decorations.markings,shapes.geometric,shapes.misc,shapes.arrows,intersections,decorations.pathmorphing,decorations.pathreplacing,backgrounds,shadows,babel}
\pgfplotsset{compat=1.18}
\usepgfplotslibrary{fillbetween}
\usepackage[np]{numprint}
\usepackage{listings,makecell,pas-tableur,twoopt}
\let\barre\undefined
\usepackage[pstricks]{tablvar}
\definecolor{darkgreen}{RGB}{0,100,0}
\providecommand{\red}{\color{red}}
\providecommand{\blue}{\color{blue}}
\providecommand{\green}{\color{green}}
\providecommand{\black}{\color{black}}
\providecommand{\white}{\color{white}}
\providecommand{\gray}{\color{gray}}
\providecommand{\cyan}{\color{cyan}}
\providecommand{\magenta}{\color{magenta}}
\providecommand{\yellow}{\color{yellow}}
\providecommand{\violet}{\color{violet}}
\providecommand{\orange}{\color{orange}}
\providecommand{\euro}{\texteuro}
\providecommand{\EUR}{\texteuro}
\providecommand{\ieme}{\textsuperscript{e}}
\providecommand{\degres}{\ensuremath{^\circ}}
\providecommand{\degree}{\ensuremath{^\circ}}
\providecommand{\vv}[1]{\overrightarrow{#1}}
\providecommand{\Coord}[1]{#1}
\providecommand{\decosix}{}
\providecommand{\starredbullet}{\ensuremath{\bullet}}
\pagestyle{empty}
\setlength{\parindent}{0pt}
\usepackage{adjustbox}\DeclareUnicodeCharacter{2236}{:}\DeclareUnicodeCharacter{2190}{\ensuremath{\leftarrow}}
\providecommand{\R}{\mathbb{R}}
\providecommand{\N}{\mathbb{N}}
\providecommand{\Z}{\mathbb{Z}}
\providecommand{\Q}{\mathbb{Q}}
\providecommand{\D}{\mathbb{D}}
\providecommand{\C}{\mathbb{C}}
\providecommand{\e}{{\mathrm{e}}}
\providecommand{\vect}[1]{\overrightarrow{#1}}
\providecommand{\vectt}[1]{\overrightarrow{\mathrm{#1}}}
\providecommand{\vvt}[1]{\overrightarrow{\mathrm{#1}}}
\providecommand{\np}[2][]{{#2~#1}}
\providecommand{\vv}[1]{{\overrightarrow{#1}}}
\renewcommand{\d}{{\mathrm{d}}}
\providecommand{\degres}{{^\circ}}
\providecommand{\degre}{{^\circ}}
\providecommand{\blue}{\color{blue}}
\providecommand{\red}{\color{red}}
\providecommand{\enskip}{\quad}
\providecommand{\llbracket}{\lbrack\!\lbrack}
\providecommand{\rrbracket}{\rbrack\!\rbrack}
\providecommand{\euro}{{\text{€}}}
\providecommand{\no}{{\text{n°}}}
\providecommand{\interval}[2]{{[#1;#2]}}
\providecommand{\intervFO}[2]{{[#1;#2[}}
\providecommand{\dotfill}{\ldots}

\providecommand{\ds}{\displaystyle}
\providecommand{\textding}[1]{#1}
\providecommand{\Oij}{\(\mathrm{O};\vec{\imath},\vec{\jmath}\)}
\providecommand{\Oijk}{\(\mathrm{O};\vec{\imath},\vec{\jmath},\vec{k}\)}
\providecommand{\Ouv}{\(\mathrm{O};\vec{u},\vec{v}\)}
\providecommand{\intervFF}[2]{\ensuremath{[#1;#2]}}
\renewcommand{\euro}{\ifmmode\text{\texteuro}\else\texteuro\fi}
\DeclareUnicodeCharacter{00B0}{\ifmmode^{\circ}\else\textdegree\fi}
\DeclareUnicodeCharacter{20AC}{\ifmmode\text{\texteuro}\else\texteuro\fi}
% Énoncés : Éducation nationale. Transcription et corrigés : APMEP, Denis Vergès et contributeurs. Découpage et classement : CoopMaths, Sébastien Lozano, Éric Elter et contributeurs. Adaptation de la présentation : Maths Explorer.
% CC BY-SA : https://coopmaths.fr/www/about/
% https://coopmaths.fr/alea/static/bac/2022/tex/bac_2022_10_sujet2_caledonie_4.tex
\begin{document}
\noindent\textit{APMEP / CoopMaths. CC BY-SA. Présentation adaptée par Maths Explorer.}\par\medskip
\textbf{Principaux domaines abordés :} suites, fonctions, primitives

\medskip

\emph{Cet exercice est un questionnaire à choix multiples.\\
Pour chacune des questions suivantes, une seule des quatre réponses proposées est exacte.\\
Une réponse fausse, une réponse multiple ou l'absence de réponse à une question ne rapporte ni n'enlève de point.\\
Pour répondre, indiquer sur la copie le numéro de la question et la lettre de la réponse choisie.\\
Aucune justification n'est demandée.}

\medskip

\begin{enumerate}
\item On considère la suite $\left(u_n\right)$ définie pour tout entier naturel $n$ par

\[u_n  = \dfrac{(- 1)^n}{n + 1}.\]

On peut affirmer que:
\begin{center}
\begin{tabularx}{\linewidth}{X X}
\textbf{a.~~} la suite $\left(u_n\right)$ diverge vers $+\infty$. &\textbf{b.~~} la suite $\left(u_n\right)$ diverge vers $-\infty$.\\
\textbf{c.~~}  la suite $\left(u_n\right)$ n'a pas de limite. &\textbf{d.~~} la suite $\left(u_n\right)$ converge.
\end{tabularx}
\end{center}

\begin{center} \decosix \decosix \decosix  \end{center}
\end{enumerate}

Dans les questions 2 et 3, on considère deux suites $\left(v_n\right)$  et $\left(w_n\right)$ vérifiant la relation : 

\[w_n = \e^{- 2v_n} + 2.\]

\begin{enumerate}[resume]
\item  Soit $a$ un nombre réel strictement positif. On a $v_0 = \ln (a)$.

\begin{center}
\begin{tabularx}{\linewidth}{X X}
\textbf{a.~~} $w_0 = \dfrac{1}{a^2}  +2$&\textbf{b.~~}$w_0 = \dfrac{1}{a^2  +2}$\\
\textbf{c.~~} $w_0 = -2a +2$&\textbf{d.~~}$w_0 = \dfrac{1}{- 2a} + 2$
\end{tabularx}
\end{center}

\item On sait que la suite $\left(v_n\right)$ est croissante. On peut affirmer que la suite $\left(w_n\right)$ est :

\begin{center}
\begin{tabularx}{\linewidth}{X X}
\textbf{a.~~}décroissante et majorée par 3.&\textbf{b.~~}décroissante et minorée par 2 .\\
\textbf{c.~~}croissante et majorée par 3 .&\textbf{d.~~}croissante et minorée par 2.
\end{tabularx}
\end{center}

\item On considère la suite $\left(a_n\right)$ ainsi définie :

\[a_0 = 2 \:\:\text{et, pour tout entier naturel}\: n, \quad a_{n+1} = \dfrac13a_n + \dfrac83.\]


Pour tout entier naturel $n$, on a :
\begin{center}
\begin{tabularx}{\linewidth}{X X}
\textbf{a.~~}$a_n = 4 \times \left(\dfrac13\right)^n - 2$&\textbf{b.~~}$a_n = - \dfrac{2}{3^n} + 4$\\
\textbf{c.~~}$a_n = 4 - \left(\dfrac13\right)^n$ & \textbf{d.~~} $a_n = 2 \times \left(\dfrac13\right)^n + \dfrac{8n}{3}$
\end{tabularx}
\end{center}
\item On considère une suite $\left(b_n\right)$ telle que, pour tout entier naturel $n$, on a :

\[b_{n+1} = b_n + \ln \left(\dfrac{2}{\left(b_n \right)^2 + 3}\right).\]

On peut affirmer que :

\begin{center}
\begin{tabularx}{\linewidth}{X X}
\textbf{a.~~} la suite $\left(b_n\right)$ est croissante.&\textbf{b.~~}la suite $\left(b_n\right)$ est décroissante.\\
\textbf{c.~~} la suite $\left(b_n\right)$ n'est pas monotone.&\textbf{d.~~}le sens de variation de la suite $\left(b_n\right)$ dépend de $b_0$.
\end{tabularx}
\end{center}

\item On considère la fonction $g$ définie sur l'intervalle $]0~;~+\infty[$ par : 

\[g(x) = \dfrac{\e^x}{x}.\]

On note $\mathcal{C}_g$ la courbe représentative de la fonction $g$ dans un repère orthogonal.

La courbe $\mathcal{C}_g$ admet :

\begin{center}
\begin{tabularx}{\linewidth}{X X}
\textbf{a.~~}une asymptote verticale
et une asymptote horizontale.&\textbf{b.~~}une asymptote verticale
et aucune asymptote horizontale.\\
\textbf{c.~~}aucune asymptote verticale et une asymptote horizontale.&\textbf{d.~~}aucune asymptote verticale et aucune asymptote horizontale.
\end{tabularx}
\end{center}
\item Soit $f$ la fonction définie sur $\R$ par

\[f(x) = x\e^{x^2+1}.\]

Soit $F$ une primitive sur $\R$ de la fonction $f$. Pour tout réel $x$, on a :

\begin{center}
\begin{tabularx}{\linewidth}{X X}
\textbf{a.~~}$F(x) = \dfrac12x^2\e^{x^2+1}$ 	& \textbf{b.~~}$F(x) = \left(1 + 2x^2 \right)\e^{x^2+1}$ \\
\textbf{c.~~}$F(x) = \e^{x^2+1}$				& \textbf{d.~~}$F(x) = \dfrac12\e^{x^2+1}$
\end{tabularx}
\end{center}
\end{enumerate}

\end{document}
